This search uses data from an integrated luminosity of 137~\ifb of \pp collisions at a centre-of-mass energy $\sqs = 13$~\TeV recorded by the ATLAS experiment during Run~2 of the LHC~(2015--2018). In this data-taking period, the number of interactions per bunch crossing ranged from about 8 to 70, with an average of 34~\cite{DAPR-2018-01}. The uncertainty in the integrated luminosity is 0.83\%~\cite{DAPR-2021-01}, obtained using the LUCID-2 detector for the primary luminosity measurements, complemented by measurements using the ID and calorimeters.
Only events recorded under stable beam conditions are used in the analysis, and general data-quality requirements ensure that all detector subsystems were operational~\cite{DAPR-2018-01}. In about 2\% of Run 2 data taking the SCT experienced a significant increase in readout errors compared to standard operation. These data periods are removed as it increases the number of spuriously reconstructed pixel-only tracks which can mimic a disappearing-track signature.

The data were collected primarily with a missing transverse momentum trigger~\cite{ATL-DAQ-PUB-2016-001,TRIG-2019-01}, to target events where the sparticle system recoils against a high-\pt ISR jet and the undetected neutral particles lead to a large amount of missing transverse momentum. The trigger threshold was raised as the luminosity increased during Run~2, but each threshold value provided an efficiency of ${>}90$\% when requiring online a reconstructed missing transverse momentum of at least 230\,GeV. Additional single-lepton ($\ell = e,\mu$) triggers are used to select regions to constrain background  estimates, and these triggers operate on an efficiency plateau when applying an offline lepton transverse momentum requirement of $\pt(\ell) > 27\,\GeV$~\cite{TRIG-2018-01,TRIG-2018-05}.

Monte Carlo (MC) simulated samples are used to model SM processes and estimate the expected signal yields. 
All samples were produced using the ATLAS simulation infrastructure~\cite{SOFT-2010-01} and either $\GEANT$~\cite{Agostinelli:2002hh} 
or a faster simulation based on a parameterisation of the calorimeter response and \GEANT for the other detector 
systems. The effect of multiple interactions in the same and neighbouring bunch crossings (\pileup) was modelled by overlaying each simulated hard-scattering event with inelastic $pp$ events generated with \PYTHIA[8.186]~\cite{Sjostrand:2007gs} using the \NNPDF[2.3lo] set of parton distribution functions (PDFs)~\cite{Ball:2012cx} and the A3 set of tuned parameters~\cite{ATL-PHYS-PUB-2016-017}. The MC simulated samples were reweighted to match the actual distribution of the mean number of interactions per bunch crossing observed in data in 2015--2018. 

The electroweakino signal samples were generated with \MGNLO[2.9.9]~\cite{Alwall:2014hca} at leading order (LO) and interfaced to \PYTHIA[8.306]~\cite{Sjostrand:2014zea} for the modelling of the parton showering (PS), hadronisation and underlying event. In the wino scenarios, $\Delta m(\chinoonepm, \ninoone)$ is dependent upon $m(\chinoonepm)$ and follows from the calculations in Refs.~\cite{Giudice:1998xp, Randall:1998uk}. In the higgsino scenarios, the mass splitting is dependent upon $m(\chinoonepm)$ and is set according to the calculations in Ref.~\cite{Nagata:2014wma}. The CMSSM $\tau$-slepton signals were generated with \MGNLO[2.9.5] at LO and interfaced to \PYTHIA[8.306], while the GMSB $\tau$-slepton signals used \MGNLO[2.6.1] at LO and \PYTHIA[8.230].

For all signal scenarios the underlying event was modelled with the A14~\cite{ATL-PHYS-PUB-2014-021} set of tuned parameters. The matrix element (ME) calculation was performed at tree level and includes the emission of up to two additional partons. The ME--PS merging was done using the CKKW-L~\cite{Lonnblad:2012ix} prescription. The \NNPDF[3.0lo] set of PDFs was used throughout~\cite{Ball:2014uwa}. To ensure efficient sample generation, a lifetime reweighting procedure was used, where a single signal-mass scenario and lifetime has a set of event weights allowing the signal to be reweighted to different proper lifetimes of the long-lived particle. The charginos and $\tau$-sleptons are propagated through the detector to their decay point using the standard simulation, including magnetic field effects and electromagnetic interactions with detector material.

Signal cross-sections are calculated to next-to-leading order (NLO) in the strong coupling constant, adding the resummation of soft gluon emission at next-to-leading-logarithm (NLL) accuracy~\cite{Beenakker:1999xh,Debove:2010kf,Fuks:2012qx,Fuks:2013vua,Fiaschi:2018hgm,Bozzi:2007qr,Fuks:2013lya,Fiaschi:2018xdm,newResummino} using \textsc{Resummino}~2.0.1. The nominal cross-section and its uncertainty are taken from an envelope of cross-section predictions using different PDF sets and factorisation and renormalisation scales, as described in Ref.~\cite{Borschensky:2014cia}. For the wino and higgsino signal samples, only processes involving the direct production of a \chinoonepm are considered in order to produce the disappearing-track signature. Therefore, the wino samples include $\chinoonep\chinoonem$ and $\chinoonepm\ninoone$ production, while the higgsino signal samples include $\chinoonep\chinoonem$, $\chinoonepm\ninoone$, and $\chinoonepm\ninotwo$. Taking an example for higgsino production with $m(\chinoonepm) = 300\,\GeV$ and $m(\ninotwo) = m(\ninoone) = 299\,\GeV$, the cross-section value for each of the $\chinoonepm\ninotwo$ and $\chinoonepm\ninoone$ processes is $90.8$\,fb, while the cross-section is $53$\,fb for $\chinoonep\chinoonem$ production.
For wino production with $m(\chinoonepm) = 300\,\GeV$ and $m(\ninoone) = 299\,\GeV$,
the cross-section values are $387$\,fb and $190$\,fb for $\ninoone\chinoonepm$ and $\chinoonep\chinoonem$ production, respectively.
Finally, the cross-section for $\tau$-slepton pair production of $m(\stau)= 300\,\GeV$ is $6.3$\,fb in both the CSSM and GMSB models.


MC simulation of $V$+\;\!jets events was used to study the fake-tracklet background. They were simulated with the \SHERPA[2.2.1]~\cite{Bothmann:2019yzt} generator using NLO matrix elements for up to two partons, and LO matrix elements for up to four partons, calculated with the Comix~\cite{Gleisberg:2008fv} and \OPENLOOPS~\cite{Buccioni:2019sur,Cascioli:2011va,Denner:2016kdg} libraries. The matrix elements were matched with the \SHERPA parton shower~\cite{Schumann:2007mg} using the \MEPSatNLO prescription~\cite{Hoeche:2011fd,Hoeche:2012yf,Catani:2001cc,Hoeche:2009rj} and the set of tuned parameters developed by the \SHERPA authors. The \NNPDF[3.0nnlo] set of PDFs~\cite{Ball:2014uwa} was used and the samples were normalised to a NNLO prediction~\cite{Anastasiou:2003ds}.

To evaluate the detector response when reconstructing short tracks and to calculate the associated reconstruction efficiencies, a set of so-called \enquote{particle gun} samples are used. These samples do not represent a specific physical event or process; instead, for each generated event the particle of interest is placed at the centre of the detector with a certain momentum and direction and can be used to investigate detector reconstruction capabilities. These samples do not include \pileup. The analysis uses three \enquote{particle gun} samples, corresponding to electrons, muons and charged pions, generated uniformly across a  transverse momentum range of $10{-}100\,\GeV$; they are used to estimate the reconstruction efficiencies as needed for the data-driven background estimation methods employed. 

